\selectlanguage{ngerman}%
\frenchspacing%
\begin{ledgroupsized}[r]{120mm}%
\footnotesize%
\pstart%
\noindent\textbf{Überlieferung:}%
\pend%
\end{ledgroupsized}%
\begin{ledgroupsized}[r]{114mm}%
\footnotesize%
\pstart%
\parindent -6mm%
\makebox[6mm][l]{\textit{F}}%
Fremdhändige Aufzeichnung mit Verbesserungen und Ergänzungen von Leibnizens Hand 
(\textit{LiF}, im edierten Text unterpunktelt):
Halle an der Saale, Universitäts- und Landesbibliothek Sachsen-Anhalt, Yg~8\textsuperscript{o}\;23~C, Bl.\;1--2
(aus der Sammlung J.\:F.\;Pfaffs).
Ein Bogen 4\textsuperscript{o},
kleine Papierabbrüche an den Falzen mit geringfügigen Textverlusten;
Wasserzeichen mittig: Hannoveraner Papier. %
Drei Seiten, einspaltig beschrieben;
Bl.\;2\,v\textsuperscript{o} ist unbeschrieben.
\pend%
\end{ledgroupsized}%
\selectlanguage{latin}%
\frenchspacing%
\vspace{8mm}%
\normalsize%
\pstart%
\noindent%
\lbrack1\,r\textsuperscript{o}\rbrack%
\phantom{Junii 1698}%
\hfill%
\phantom{\lbrack1\,r\textsuperscript{o}\rbrack}%
\edtext{%
\dotuline{Junii 1698}%
}{%
\lemma{\dotuline{Junii 1698}}%
\Bfootnote{%
\textit{erg.\;LiF}%
}}
\pend%
\pstart%
\noindent%
In
\edtext{%
duobus
\dotuline{Tubis}%
\indexS{tubus}%
}{%
\lemma{duobus}%
\Bfootnote{%
\textit{(1)}~barometris\indexS{barometrum}%
~\textit{F} %
\textit{(2)}~\dotuline{Tubis}%
~\textit{LiF}%
}}
\textit{AC} et \textit{BD} majori et minori,
tentamen%
\indexS{tentamen} institutum fuit
\edtext{%
quousque quantitate aliqua\;%
\edtext{%
\mercury ii%
\indexS{mercurius in tubo}%
}{%
\lemma{\mercury ii}%
\Cfootnote{%
Das Symbol \mercury\ steht für \textit{mercur}(\textit{ius}).%
}}
compressus aer%
\indexS{aer compressus in tubo}%
}{%
\lemma{quousque}%
\Bfootnote{%
\textit{(1)}~certa %
\textit{(2)}~quantitate aliqua \mercury ii %
\bars{\textlangle indefinita\textrangle\ \textit{gestr.\;LiF}} %
compressus aer%
~\textit{F}%
}}
per aequalem et determinatum caloris gradum%
\indexS{gradus caloris}%
\indexS{calor}
se extendere%
\indexS{aer extendens se per calorem}%
\indexS{extensio aeris per calorem}
\mercury iumque%
\indexS{mercurius in tubo}
attollere possit.
Processus fuit talis.%
\indexS{processus tentaminis}%
\pend%
\pstart%
\edtext{%
\edtext{%
In \dotuline{Tubo}%
\indexS{tubus}%
}{%
\lemma{In}%
\Bfootnote{%
\textit{(1)}~barometro\indexS{barometrum}%
~\textit{F} %
\textit{(2)}~\dotuline{Tubo}%
~\textit{LiF}%
}}
\textit{AC},%
}{%
\lemma{In \dotuline{Tubo} \textit{AC}}%
\Cfootnote{%
Die Beschreibung des Experiments bezieht sich auf ein in N.\;\ref{36650} nicht überliefertes Diagramm.%
}}
spatium aeris in globulo \textit{C}%
\indexS{globulus tubi}
comprimendi \textit{pq}%
\indexS{aer comprimendus in tubo}
continebat
1 \Pfund,%
\indexS{Pfund}
20
\edtext{%
loth,%
\indexS{Loth}%
}{%
\lemma{loth}%
\Cfootnote{%
Gewichtseinheit, die gemäß dem Handelsgewicht $\displaystyle\frac{1}{32}$ des Pfundes entsprach.%
}}
$1 \frac{1}{2}$~%
\edtext{%
qt.%
\indexS{Quentchen}%
}{%
\lemma{qt}%
\Cfootnote{%
Quentchen, d.\,h. Gewichtseinheit, die gemäß dem Handelsgewicht einem Viertel %
des Lots bzw. $\displaystyle\frac{1}{128}$ des Pfundes entsprach.
Hierfür wird im Folgenden zumeist \textit{qtchen} als Abkürzung verwendet.%
}}%
\;\mercury ii.
In altero minori \textit{BD},%
\indexS{tubus}
spatium \textit{lm} globuli \textit{D}%
\indexS{globulus tubi}
replebatur
1 \Pfund,%
\indexS{Pfund}
7 loth,%
\indexS{Loth}
$2 \frac{1}{2}$ qtchen%
\indexS{Quentchen}%
\;\mercury ii.%
\indexS{mercurius in tubo}%
\pend%
\pstart%
Hisce exemtis ponderatisque
infundebatur de
\edtext{%
novo in \dotuline{Tubum}%
\indexS{tubus}%
}{%
\lemma{novo in}%
\Bfootnote{%
\textit{(1)}~barometrum\indexS{barometrum}%
~\textit{F} %
\textit{(2)}~\dotuline{Tubum}%
~\textit{LiF}%
}}
\textit{AC}
\edtext{%
\mercury \dotuline{ius},%
\indexS{mercurius in tubo}%
\indexS{mercurius infusus}%
}{%
\lemma{\dotuline{ius},}%
\Bfootnote{%
\textit{erg.\;LiF}%
}}
donec in tubulo \textit{A}%
\indexS{tubulus}
in \textit{i},
in globulo \textit{C}%
\indexS{globulus tubi}
vero in \textit{r}
u\textlangle sque\textrangle\
ascendisset.%
\pend%
\pstart%
Sic et in altero \textit{BD},
\edtext{%
\mercury \dotuline{ius}%
\indexS{mercurius in tubo}%
}{%
\lemma{\dotuline{ius}}%
\Bfootnote{%
\textit{erg.\;LiF}%
}}
infusus%
\indexS{mercurius infusus}
in tubulo \textit{B}%
\indexS{tubulus}
locum \textit{g},
in globulo \textit{D}%
\indexS{globulus tubi}
vero locum \textit{o}
\edlabel{Halle_ULB_Yg_8o_23_C_001-002_001r_wlw-1}%
attingebat.%
\edtext{}{%
\xxref{Halle_ULB_Yg_8o_23_C_001-002_001r_wlw-1}{Halle_ULB_Yg_8o_23_C_001-002_001r_wlw-2}%
\lemma{attingebat.}%
\Bfootnote{%
\lbrack/\rbrack\ %
\textit{(1)}~Aer%
~\textit{F} %
\textit{(2)}~\dotuline{Aeris}%
~\textit{LiF} %
igitur in
\bars{barometro \textit{ändert Hrsg.}}
\textit{AC} %
\bars{in locum \textit{rq} \textit{streicht Hrsg.}} %
\textit{(a)}~compressus%
~\textit{F} %
\textit{(b)}~\dotuline{compressi spatium \textit{rq}}%
~\textit{LiF} %
post %
\bars{totam \textit{gestr.\;LiF}} %
operationem %
\textit{(aa)}~mensuratus,%
~\textit{F} %
\textit{(bb)}~\dotuline{mensuratum},%
~\textit{LiF}%
}}%
\pend%
\newpage%
\pstart%
\dotuline{Aeris} igitur%
\indexS{aer compressus in tubo}
in \lbrack tubo\rbrack\
\textit{AC}%
\indexS{barometrum}
\dotuline{compressi%
\indexS{aer compressus in tubo}
spatium \textit{rq}}
post operationem%
\indexS{operatio in tentamine}
\dotuline{mensuratum,}%
\edlabel{Halle_ULB_Yg_8o_23_C_001-002_001r_wlw-2}
1 \Pfund,%
\indexS{Pfund}
5 loth,%
\indexS{Loth}
$2 \frac{1}{2}$ qtchen%
\indexS{Quentchen}
\mercury ii
capiebat,
\edtext{%
cum
\dotuline{alterius}%
\indexS{aer adimplens spatium}
in \textit{BD}
spatium \textit{om}
\dotuline{adimplentis capacitas}%
\indexS{capacitas aeris compressi}
non nisi
\mercury ium%
\indexS{mercurius in tubo}%
\indexS{mercurius infusus}
$29 \frac{1}{2}$ loth%
\indexS{Loth}
\dotuline{ponderantem caperet.}%
}{%
\lemma{cum}%
\Bfootnote{%
\textit{(1)}~alter%
~\textit{F} %
\textit{(2)}~\dotuline{alterius}%
~\textit{LiF} %
in %
\bars{barometro\indexS{barometrum} \textit{gestr.\;LiF}} %
\textit{BD} spatium \textit{om} %
\textit{(a)}~adimplens%
~\textit{F} %
\textit{(b)}~\dotuline{adimplentis capacitas}%
~\textit{LiF} %
non nisi \mercury ium $29 \frac{1}{2}$ loth %
\textit{(aa)}~ponderare.%
~\textit{F} %
\textit{(bb)}~\dotuline{ponderantem caperet}.%
~\textit{LiF}%
}}%
\pend%
\pstart%
Immittebantur nunc ambo in
\lbrack1\,v\textsuperscript{o}\rbrack\ %
\edtext{%
$\bigtriangledown$m%
}{%
\lemma{$\bigtriangledown$m}%
\Cfootnote{%
Das Symbol $\bigtriangledown$ steht für \textit{aqua}.%
}}
calidam,%
\indexS{aqua calida}
ascendebatque
\edtext{%
\mercury%
\indexS{mercurius in tubo}%
\indexS{mercurius ascendens calore}
in
\dotuline{Tubo}%
\indexS{tubus}
\textit{AC} in~\textit{k},
in
\dotuline{Tubo}%
\indexS{tubus}%
\textit{BD} in \textit{h},%
}{%
\lemma{\mercury in}%
\Bfootnote{%
\textit{(1)}~barometro\indexS{barometrum}%
~\textit{F} %
\textit{(2)}~\dotuline{Tubo}%
~\textit{LiF} %
\textit{AC} in~\textit{k}, in %
\textit{(a)}~minori%
~\textit{F} %
\textit{(b)}~\dotuline{Tubo}%
~\textit{LiF} %
\textit{BD} in \textit{h},%
~\textit{F}%
}}
usque Spatium
\edtext{%
\textit{ik}
\dotuline{capiebat
\mercury ii}%
}{%
\lemma{\textit{ik}}%
\Bfootnote{%
\textit{(1)}~in majori%
~\textit{F} %
\textit{(2)}~\dotuline{capiebat \mercury ii}%
~\textit{LiF}%
}}
6 loth%
\indexS{Loth}
$3 \frac{1}{2}$ qtchen,%
\indexS{Quentchen}
\textit{gh}
\edtext{%
vero
\dotuline{Mercurii}%
\indexS{mercurius in tubo}%
\indexS{mercurius ascendens calore}%
}{%
\lemma{vero}%
\Bfootnote{%
\textit{(1)}~in minori%
~\textit{F} %
\textit{(2)}~\dotuline{ca} %
\textit{(3)}~\dotuline{\mercury ii} %
\textit{(4)}~\dotuline{Mercurii}%
~\textit{LiF}%
}}
$3 \frac{1}{2}$~loth%
\indexS{Loth}
capiebat.%
\pend%
\pstart%
\edtext{%
Eaedem
\dotuline{hae capacitates,}%
\indexS{capacitas aeris compressi}%
}{%
\lemma{Eadem}%
\Bfootnote{%
\textit{(1)}~haec pondera,\indexS{pondus mercurii infusi}%
~\textit{F} %
\textit{(2)}~\dotuline{hae capacitates,}%
~\textit{LiF}%
}}
nimirum
\edtext{%
\dotuline{pro}%
}{%
\lemma{\dotuline{pro}}%
\Bfootnote{%
\textit{erg.\;LiF}%
}}
6 loth%
\indexS{Loth}
$3 \frac{1}{2}$ qtchen%
\indexS{Quentchen}
et
\edtext{%
\dotuline{pro}%
}{%
\lemma{\dotuline{pro}}%
\Bfootnote{%
\textit{erg.\;LiF}%
}}
$3 \frac{1}{2}$ loth,%
\indexS{Loth}
addendae
\edtext{%
sunt
\dotuline{capacitatibus%
\indexS{capacitas aeris compressi}
seu spatiis}%
}{%
\lemma{sunt}%
\Bfootnote{%
\textit{(1)}~ponderibus %
~\textit{F} %
\textit{(2)}~\dotuline{capacitatibus seu spatiis}%
~\textit{LiF}%
}}
compressi aeris \textit{rq}%
\indexS{aer compressus in tubo}
capienti %
1 \Pfund,%
\indexS{Pfund}
5 loth,%
\indexS{Loth}
$2 \frac{1}{2}$
\edtext{%
qtchen%
\indexS{Quentchen}
\dotuline{in}%
}{%
\lemma{qtchen}%
\Bfootnote{%
\textit{(1)}~in barometro\indexS{barometrum}%
~\textit{F} %
\textit{(2)}~\dotuline{in}%
~\textit{LiF}%
}}
\textit{AC},
et \textit{om}
capienti %
$\langle2\rangle9 \frac{1}{2}$
\edtext{%
loth%
\indexS{Loth}
in \textit{BD}.%
}{%
\lemma{loth in}%
\Bfootnote{%
\bars{barometro \textit{gestr.\;LiF}}
\textit{BD}.%
~\textit{F}%
}}
Prodibit in illo
\edtext{%
\dotuline{capacitas}%
\indexS{capacitas aeris compressi}%
}{%
\lemma{\dotuline{capacitas}}%
\Bfootnote{%
\textit{erg.\;LiF}%
}}
1 \Pfund,%
\indexS{Pfund}
12 loth,%
\indexS{Loth}
2 qtchen,%
\indexS{Quentchen}
in hoc
\edtext{%
\dotuline{capacitas pro}%
}{%
\lemma{\dotuline{capacitas pro}}%
\Bfootnote{%
\textit{erg.\;LiF}%
}}
1 \Pfund,%
\indexS{Pfund}
1 loth.%
\indexS{Loth}
Quia
quantum in ambobus assurgit
\mercury,%
\indexS{mercurius assurgens calore}
tantum dilatatur utrobique spatium
et per consequens aer%
\indexS{aer dilatatus calore}
in globulis \textit{C} et \textit{D}.%
\indexS{globulus tubi}
\pend%
\pstart%
Hisce
\edtext{%
peractis,
in \textit{BD}%
}{%
\lemma{peractis, in}%
\Bfootnote{%
\textit{(1)}~minus barometrum\indexS{barometrum}%
~\textit{F} %
\textit{(2)}~\bars{\dotuline{instrume}\indexS{instrumentum} \textit{erg.\;u. gestr.\;LiF}} %
\textit{BD}%
~\textit{F}%
}}
adhuc major quantitas
\mercury ii%
\indexS{mercurius infusus}
infundebatur,
ita ut in tubulo \textit{B}%
\indexS{tubulus}
locum \textit{e},
in globulo \textit{D}%
\indexS{globulus tubi}
vero locum \textit{n}
attingeret.
Notabantur ambo exterius.
\lbrack2\,r\textsuperscript{o}\rbrack\ %
\pend%
\pstart%
Imittebatur iterum in
$\bigtriangledown$m%
\indexS{aqua calida}
eundem caloris gradum,%
\indexS{gradus caloris}%
\indexS{calor}
beneficio thermometri%
\indexS{thermometrum}
determinatum,
cum pristina habentem,
ascendebatque in \textit{f} usque
\mercury ius.%
\indexS{mercurius ascendens calore}%
\indexS{mercurius in tubo}%
\pend%
\pstart%
Mensuratum
\mercury ii
spatium
\edtext{%
\textit{ef}
$\dotuline{2} \frac{1}{2}$\dotuline{ loth,}%
\indexS{Loth}%
}{%
\lemma{\textit{ef}}%
\Bfootnote{%
\textit{(1)}~non nisi%
~\textit{F} %
\textit{(2)}~$\dotuline{2} \frac{1}{2}$\dotuline{ loth,}%
~\textit{LiF}%
}}
spatium vero \textit{nm}
aeris postremo
\edtext{%
compres\lbrack s\rbrack i%
\indexS{aer compressus in tubo}%
}{%
\lemma{compresi}%
\Bfootnote{%
\textit{F\;ändert Hrsg.}%
}}
19~loth%
\indexS{Loth}
1 qtchen%
\indexS{Quentchen}
continebat.%
\pend%
\pstart%
Illud huic additum
producebat iterum totum spatium
quod calore dilatatus%
\indexS{aer dilatatus calore}
aer occupabat%
\indexS{aer compressus in tubo}%
\lbrack:\rbrack\
21 loth%
\indexS{Loth}
3 qtchen.%
\indexS{Quentchen}
\pend%
\pstart%
Tandem et spatia
quae \mercury%
\indexS{mercurius in tubo}
a regione dimidiorum hemisphaeriorum \textit{S}, \textit{S}%
\indexS{hemisphaerium}
usque in \textit{i}, \textit{g} et \textit{e} occupaverat,
in scala pedali%
\indexS{scala pedalis}
\edtext{%
e\textlangle m\textrangle etiebantur,
inveniebaturque%
}{%
\lemma{e\textlangle m\textrangle etiebantur,}%
\Bfootnote{%
\textit{(1)}~eratque %
\textit{(2)}~inveniebaturque%
~\textit{F}%
}}
\textit{Si} 4 ped.,%
\indexS{pes}
\textit{Sg} 1 ped.
$6 \frac{7}{10}$ digit.,%
\indexS{digitus}
\textit{Se} 4 ped.%
\indexS{pes}
$1 \frac{1}{2}$ digit.
Spatia vero
\textit{ik} $4 \frac{1}{7}$ digit.,
\textit{gh} 3 digit.%
\indexS{digitus}
et \textit{ef} $2 \frac{1}{10}$ erant.%
\pend%
\newpage%
\pstart%
\edtext{%
\dotuline{%
Aer per \textit{Si}
compressus est%
\indexS{aer compressus in tubo}
in spatium \textit{qr}
quod est ad spatium \textit{qp}
aeris ejusdem naturalem statum habentis,%
\indexS{aer habens statum naturalem}%
\indexS{status aeris naturalis}
ut \mercury ii quantitates%
\indexS{mercurius in tubo}%
\indexS{mercurius contentus ab aere compresso}
quas continent haec spatia,
seu ut 1 \Pfund%
\indexS{Pfund}
5~loth%
\indexS{Loth}
}%
$\dotuline{2}
\frac{1}{2}$%
\dotuline{ quentchen%
\indexS{Quentchen}
ad 1 \Pfund%
\indexS{Pfund}
20 loth%
\indexS{Loth}
}%
$\dotuline{1}
\frac{1}{2}$%
\dotuline{\ quentchen.%
\indexS{Quentchen}
Dividatur hoc per illud,
fit%
}%
}{%
\lemma{\dotuline{Aer per \textit{Si} compressus}}%
\Bfootnote{%
\textit{(1)}~\dotuline{\textit{S}}\textlangle\dotuline{--}\textrangle\ %
\textit{(2)}~\dotuline{est in} %
\textit{(a)}~\dotuline{spatia} %
\textit{(b)}~\dotuline{spatium \textit{qr} quod est} %
\bars{\dotuline{ad} \textit{erg.}} \dotuline{spatium \textit{qp}} %
\textit{(aa)}~\dotuline{natural} %
\textit{(bb)}~\dotuline{aeris ejusdem naturalem statum habentis,} %
\textit{(aaa)}~\dotuline{ut Mercurius quem} %
\textit{(bbb)}~\dotuline{ut \mercury ii} %
\lbrack...\rbrack\ \dotuline{5 loth }$\dotuline{2} \frac{1}{2}$ %
\textit{(aaaa)}~\dotuline{q}\textlangle\dotuline{\textendash}\textrangle\ %
\textit{(bbbb)}~\dotuline{quentchen ad} %
\lbrack...\rbrack\ \dotuline{illud, fit} %
\textit{erg.\;LiF}%
}}%
\lbrack:\rbrack\
\edtext{%
$1 \frac{118}{301}$
sive
$\frac{419}{301}$.%
}{%
\lemma{$1 \frac{118}{301}$}%
\Bfootnote{%
\bars{L\,qtchen \textit{gestr.\;LiF}} %
sive $\frac{419}{301}$ %
\bars{L\,qtchen \textit{gestr.\;LiF}}.%
~\textit{F}%
}}%
\pend%
\pstart%
\edtext{%
\dotuline{%
Aer per \textit{Sg}
compressus est%
\indexS{aer compressus in tubo}
in spatium \textit{mo}
quod est ad spatium \textit{ml}
aeris ejusdem naturalis statum habentis,%
\indexS{aer habens statum naturalem}%
\indexS{status aeris naturalis}
ut \mercury ii quantitates%
\indexS{mercurius in tubo}%
\indexS{mercurius contentus ab aere compresso}
quas continent haec spatia,
seu ut }%
$\dotuline{29}
\frac{1}{2}$%
\dotuline{ loth%
\indexS{Loth}
ad 1 \Pfund%
\indexS{Pfund}
7 loth%
\indexS{Loth}
}%
$\dotuline{2}
\frac{1}{2}$%
\dotuline{ quentchen.%
\indexS{Quentchen}
Dividatur hoc per illud,
fit}%
}{%
\lemma{\dotuline{Aer per}}%
\Bfootnote{%
\lbrack...\rbrack\ \dotuline{illud, fit} %
\textit{erg.\;LiF}%
}}%
\lbrack:\rbrack\
\edtext{%
$1 \frac{81}{236}$
sive
$\frac{317}{236}$.%
}{%
\lemma{$1 \frac{81}{236}$}%
\Bfootnote{%
\bars{L\,qtchen \textit{gestr.\;LiF}} %
sive
$\frac{317}{236}$ %
\bars{L\,qtchen \textit{gestr.\;LiF}}.%
~\textit{F}%
}}%
\pend%
\pstart%
\edtext{%
\dotuline{%
Aer per \textit{Se}
compressus est%
\indexS{aer compressus in tubo}
in spatium \textit{mn}
quod est ad spatium \textit{ml}
aeris ejusdem naturalis statum habentis,%
\indexS{aer habens statum naturalem}%
\indexS{status aeris naturalis}
ut \mercury ii quantitates%
\indexS{mercurius in tubo}%
\indexS{mercurius contentus ab aere compresso}
quas continent haec spatia,
seu ut 19 loth%
\indexS{Loth}
1 quentchen%
\indexS{Quentchen}
ad 1 \Pfund%
\indexS{Pfund}
7 loth%
\indexS{Loth}
}%
$\dotuline{2} \frac{1}{2}$%
\dotuline{ quentchen.%
\indexS{Quentchen}
Dividatur hoc per illud,
fit%
}%
}{%
\lemma{\dotuline{Aer per}}%
\Bfootnote{%
\lbrack...\rbrack\ \dotuline{haec spatia,} %
\textit{(1)}~\dotuline{seu ut 29 loth 1 quentchen ad 1 \Pfund\ 7 l,} %
\textit{(2)}~\dotuline{seu ut 19} %
\lbrack...\rbrack\ \dotuline{illud, fit} %
\textit{erg.\;LiF}%
}}%
\lbrack:\rbrack\
\edtext{%
$2 \frac{9}{154}$
sive
$\frac{317}{154}$.%
}{%
\lemma{$2 \frac{9}{154}$}%
\Bfootnote{%
\bars{L\,qtchen \textit{gestr.\;LiF}} sive $\frac{317}{154}$ %
\bars{L\,qtchen \textit{gestr.\;LiF}}.%
~\textit{F}%
}}%
\pend%
\selectlanguage{latin}%
\frenchspacing%
\normalsize%
